\chapter{
  Statement on effectiveness of your own performance and the team performance
}
A large portion of this project required independent work, particularly in the
integration of COMSOL Multiphysics with a genetic algorithm written in Java. I
was solely responsible for designing and executing the simulation workflow,
which involved building a programmatic interface using the COMSOL Java API,
developing and tuning the genetic algorithm, and handling large volumes of
simulation data. This independence meant I had to manage my time efficiently,
troubleshoot errors without immediate supervision, and learn new technical
skills on my own, including how to script mesh-free geometry updates, automate
parametric sweeps, and extract electromagnetic field data for postprocessing.
This project gave me the opportunity to demonstrate and refine my
problem-solving skills. I had to carefully isolate variables, run controlled
tests, and develop practical workarounds. This trial-and-error approach deepened
my understanding of both the simulation tools and the underlying physics.
Although the technical work was done mostly independently, collaboration was
still important. Early discussions with post-graduates helped me get set up with
the software and avoid common pitfalls. I benefited especially from informal
peer discussions which helped me learn valuable debugging strategies, simulation
efficiency tricks. Weekly supervisor meetings in the first term were especially
useful for developing the foundational understandings of MASERs and helped keep
the project on track.
